\documentclass[conference]{IEEEtran}
\IEEEoverridecommandlockouts

\usepackage{cite}
\usepackage{amsmath,amssymb,amsfonts}
\usepackage{algorithmic}
\usepackage{graphicx}
\usepackage{textcomp}
\usepackage{xcolor}
\usepackage{booktabs}
\usepackage{multirow}
\usepackage{bm}

\def\BibTeX{{\rm B\kern-.05em{\sc i\kern-.025em b}\kern-.08em
    T\kern-.1667em\lower.7ex\hbox{E}\kern-.125emX}}

\begin{document}

\title{DMRL: Polarity-Coupled Optimal Transport with\\
Disagreement-Aware Fusion for Multimodal Sentiment Analysis}

\author{\IEEEauthorblockN{Anonymous Authors}
\IEEEauthorblockA{\textit{Anonymous Institution} \\
Anonymous City, Country \\
anonymous@email.com}}

\maketitle

\begin{abstract}
Multimodal Sentiment Analysis (MSA) aims to infer human sentiment by integrating
textual, acoustic, and visual signals. Non-verbal cues may either corroborate or
contradict the text, so the central challenge is to align cross-modal evidence
while characterizing whether aligned cues agree or conflict. Existing methods
typically address these with two separate mechanisms: they split representations
into shared/private (or consistent/complementary) parts via subspace decomposition
or contrastive learning, and then fuse them with attention or routing. This
introduces many components to co-tune and decouples ``alignment'' from
``consensus/conflict judgment.'' We propose \textbf{DMRL}, a framework whose sole
core is \emph{Polarity-Coupled Optimal Transport} (PCOT). DMRL first extracts a few
sparse evidence slots per modality, each carrying only a \emph{polarity} and a
\emph{reliability}. PCOT then couples a semantic distance and a sentiment-polarity
distance into a single transport cost and solves a soft cross-modal matching with
entropic Sinkhorn. Crucially, the \emph{same} transport plan is split, by polarity
consistency, into a \emph{consensus flow} and a \emph{conflict flow}; the relative
mass of the two flows directly yields a closed-form, sample-level \emph{cross-modal
disagreement score}. This score in turn gates a \emph{disagreement-aware fusion} of
the consensus and conflict streams, and a compact ordinal distribution head
produces the sentiment. Cross-modal alignment, consensus/conflict decomposition,
and disagreement measurement are thus unified in a single transport, with no extra
decomposer or router. Experiments on CMU-MOSI, CMU-MOSEI, and CH-SIMS show that
DMRL substantially simplifies the architecture while matching or surpassing strong
baselines.
\end{abstract}

\begin{IEEEkeywords}
multimodal sentiment analysis, optimal transport, polarity coupling,
consensus--conflict decomposition, disagreement-aware fusion
\end{IEEEkeywords}

\section{Introduction}
Multimodal Sentiment Analysis (MSA) seeks to understand human emotions by jointly
modeling language, acoustic, and visual streams that co-occur in conversational
video~\cite{zadeh2018mosei,yu2020ch}. Because non-verbal cues frequently reinforce,
weaken, or even reverse the apparent polarity of spoken words, effectively aligning
and fusing cross-modal evidence is central to the task.

A dominant paradigm is to ``decompose, then fuse'': via shared--private subspace
decomposition~\cite{hazarika2020misa}, contrastive learning, or information
bottlenecks, each modality's representation is split into a cross-modal consistent
part and a modality-specific (or complementary/conflicting) part, which are then
combined by attention or dynamic routing. Despite strong progress, two issues
remain. First, \emph{alignment and consensus/conflict judgment are decoupled}.
``Which cue aligns with which'' is usually decided by one similarity or attention
mechanism, whereas ``whether an aligned pair is consensual or conflicting'' is
judged by another subspace or contrastive objective. Being independent, they add
components that must be co-tuned and leave the notion of consensus versus conflict
without a unified, interpretable origin. Second, \emph{the degree of disagreement
is hard to measure explicitly}. How inconsistent the text and the non-verbal cues
actually are in a given sample is usually only implicit in the fusion weights,
lacking a direct, closed-form, sample-level disagreement measure to guide
fusion---yet it is precisely this disagreement that determines whether the model
should trust consensual signals or confront conflicting ones.

We argue that a more faithful approach should unify \emph{alignment} and
\emph{consensus/conflict decomposition} within a single cross-modal matching, and
read the disagreement strength directly from that matching to guide fusion. To this
end, we propose \textbf{DMRL}, whose sole core is \emph{Polarity-Coupled Optimal
Transport} (PCOT).

DMRL first maps modalities into a shared space through modality-specific
multi-scale stems, a shared Transformer encoder, and lightweight adapters, and
extracts a few \emph{sparse evidence slots} per modality, each carrying only two
interpretable scalar attributes: a \emph{polarity} $p\in[-1,1]$ and a
\emph{reliability} $r\ge 0$. PCOT then uses a ``semantic distance + polarity
distance'' transport cost and Sinkhorn to obtain a soft matching plan $\mathbf{T}$;
a polarity-consistency gate splits the \emph{same} $\mathbf{T}$ into a
\emph{consensus flow} $\mathbf{T}^{+}$ that aggregates same-polarity aligned
evidence and a \emph{conflict flow} $\mathbf{T}^{-}$ that characterizes
opposite-polarity aligned evidence. The relative mass carried by the two flows
naturally gives a closed-form, sample-level \emph{disagreement score}
$\delta\in[0,1]$. Finally, $\delta$ gates a \emph{disagreement-aware fusion} of the
consensus and conflict representations, and an \emph{ordinal distribution head}
over fixed sentiment anchors produces the score.

Our contributions are summarized as follows:
\begin{itemize}
\item We introduce \emph{Polarity-Coupled Optimal Transport} (PCOT), which couples
semantic similarity and sentiment-polarity agreement into a single transport cost
and splits the resulting plan into consensus and conflict flows, thus performing
cross-modal alignment and consensus/conflict decomposition with \emph{one}
mechanism, without a separate subspace decomposer or contrastive branch.
\item We derive a \emph{closed-form, sample-level cross-modal disagreement score}
from the relative mass of the conflict flow, and use it to gate a
disagreement-aware fusion of the consensus and conflict streams; the score needs no
extra routing network and serves as an interpretable diagnostic.
\item We give a compact instantiation (polarity--reliability evidence slots, an
ordinal distribution head, and a polarity-grounding constraint) that removes
multi-branch complexity while preserving accuracy, and is competitive with or
superior to strong baselines on three standard benchmarks.
\end{itemize}

\section{Related Work}

\subsection{Multimodal Sentiment Analysis and Representation Decomposition}
Early MSA research modeled cross-modal interaction through tensor
fusion~\cite{zadeh2017tensor} and cross-modal attention~\cite{tsai2019mult}. To
control redundancy and preserve discriminative cues, a prominent line decomposes
representations into cross-modal shared and modality-specific subspaces, exemplified
by MISA~\cite{hazarika2020misa}; subsequent work refines this decomposition:
FINE~\cite{liu2025fine} factorizes shared and private representations via mutual
information estimation and suppresses task-irrelevant noise, TSD~\cite{meng2026tsd}
introduces pairwise-shared subspaces beyond global-shared and private ones to model
pairwise synergy, DLF~\cite{wang2025dlf} disentangles shared/specific information
around a language-centered design with geometric constraints, and
DERL~\cite{li2026derl} uses disentangled experts for robustness under missing
modalities. Another line explicitly models ``consensus versus
conflict/complementarity'': DCAF~\cite{lv2026dcaf} separates similar and
complementary representations via contrastive decoupling and adaptively fuses them,
DCCF~\cite{zhou2025dccf} decouples consistency and complementarity information to
alleviate their optimization conflict, and PRISM~\cite{su2026prism} models
complementarity, conflict, and reliability differences among evidence in a shared
prototype space with dynamic re-weighting.

These methods mostly organize evidence along a \emph{representation axis}
(shared/private or consistent/complementary), and ``alignment (who matches whom)''
and ``decomposition (which class the match belongs to)'' are handled by different
mechanisms. DMRL differs: it organizes evidence along \emph{polarity consistency}
and unifies alignment and consensus/conflict decomposition within a single optimal
transport---consensus and conflict are not produced by an extra subspace or
contrastive decoder but are the direct result of splitting the \emph{same} transport
plan by a polarity gate; it also keeps a few decisive evidence slots carrying
polarity and reliability, rather than a single per-modality vector or prototype.

\subsection{Optimal Transport for Cross-Modal Alignment}
Optimal transport (OT) and its entropic Sinkhorn approximation~\cite{cuturi2013sinkhorn}
provide a principled tool for soft matching between two sets of elements. In MSA,
PaSE~\cite{he2026pase} aligns prototype representations with entropic OT and
combines Shapley-based gradient modulation to mitigate modality competition, a
representative recent use of OT in MSA. DMRL's transport differs in two essential
ways: first, the transport cost includes not only a semantic distance but also
\emph{couples a sentiment-polarity distance}, so that matching naturally
distinguishes same- from opposite-polarity evidence; second, we do not stop at an
alignment plan but further split the \emph{same} plan by polarity consistency into
consensus and conflict flows, reading a disagreement score from the ratio of the
two flow masses. To our knowledge, coupling polarity into the transport cost and
deriving consensus/conflict representations and a disagreement measure directly from
the mass distribution of the transport plan has not been reported before.

\subsection{Uncertainty and Ordinal Modeling}
Trustworthy multimodal learning estimates predictive uncertainty for
reliability-aware fusion, e.g., through evidential or mixture-of-Gaussian
formulations~\cite{han2021trusted}. In MSA, TMSON~\cite{xie2024tmson} estimates a
per-modality uncertainty distribution and fuses via Bayes' rule while constraining
prediction in an ordinal sentiment space. DMRL retains ordinal modeling as a simple
yet effective ingredient: it directly predicts a distribution over fixed sentiment
anchors and uses its expectation as the score; but unlike methods that rely on an
explicit uncertainty distribution to weight fusion, DMRL uses the \emph{disagreement
score} read from the transport plan to modulate the fusion of consensus and conflict
representations, in a more compact form.

\section{Method}
Figure~\ref{fig:overview} (not shown) summarizes DMRL. Given a text sequence
$\mathbf{X}_t$, an acoustic sequence $\mathbf{X}_a$, and a visual sequence
$\mathbf{X}_v$, DMRL proceeds through (i) shared encoding with modality adapters,
(ii) sparse polar evidence extraction, (iii) polarity-coupled optimal transport
with consensus/conflict flow decomposition, and (iv) disagreement-aware fusion and
ordinal distribution prediction.

\subsection{Shared Encoding with Modality Adapters}
For text, a pretrained BERT/RoBERTa encoder produces token features. Each modality
$m \in \{t,a,v\}$ is first passed through a multi-scale temporal stem composed of
parallel $1\times1$, $3\times3$, and $5\times5$ depthwise-style 1D convolutions
with a residual projection, mapping all modalities into a common $d$-dimensional
space:
\begin{equation}
\mathbf{H}_m^{0} = \mathrm{Stem}_m(\mathbf{X}_m) + \mathbf{e}_m,
\end{equation}
where $\mathbf{e}_m$ is a learnable modality embedding. A \emph{shared}
mask-aware Transformer encoder $\mathcal{E}$ then processes all modalities, which
reduces parameters and encourages a common evidence space, followed by a
lightweight modality-specific residual adapter $\mathcal{A}_m$:
\begin{equation}
\mathbf{H}_m = \mathcal{A}_m\big(\mathcal{E}(\mathbf{H}_m^{0})\big).
\end{equation}

\subsection{Sparse Polar Evidence Extraction}
From $\mathbf{H}_m \in \mathbb{R}^{T_m \times d}$ we extract $K$ evidence slots via
slot attention with \emph{hybrid static--adaptive queries}. Static slot queries
$\mathbf{Q}^{0}\in\mathbb{R}^{K\times d}$ are modulated by the global context
$\mathbf{c}_m = \mathrm{MaskedMean}(\mathbf{H}_m)$ through a gating function:
\begin{align}
\mathbf{Q} &= \mathrm{LN}\!\left(\mathbf{Q}^{0} + \bm{g}\odot \phi(\mathbf{c}_m)\right),\\
\bm{g} &= \sigma\!\left(\mathrm{MLP}([\mathbf{Q}^{0};\phi(\mathbf{c}_m)])\right).
\end{align}
After masked softmax we obtain $\mathbf{A}\in\mathbb{R}^{K\times T_m}$, and the
evidence vectors are $\mathbf{Z} = \mathbf{A}\,\mathbf{V}(\mathbf{H}_m)$. Each slot
$i$ carries only \emph{two} interpretable scalar attributes:
\begin{equation}
p_i = \tanh(\cdot)\in[-1,1]\ \text{(polarity)}, \qquad
r_i = \mathrm{softplus}(\cdot)\ge 0\ \text{(reliability)}.
\end{equation}
The reliability yields pooling weights $\bm{w} = \mathrm{softmax}(\bm{r})$, so that
the modality summary is $\bar{\mathbf{z}}_m = \sum_i w_i \mathbf{z}_i$. Compared with
attaching strength, uncertainty, and temporal position to every slot, we keep only
the polarity that drives transport and the reliability that drives pooling, making
the evidence representation markedly more compact.

\subsection{Polarity-Coupled Optimal Transport}
PCOT is the core of DMRL. For an ordered modality pair $(a,b)$ (we anchor on text,
i.e., $(t,a)$ and $(t,v)$), let the semantic similarity of normalized evidence be
$\mathrm{sim}_{ij} = \hat{\mathbf{z}}^{a}_i \cdot \hat{\mathbf{z}}^{b}_j$. The
transport cost \emph{couples} a semantic distance and a sentiment-polarity distance:
\begin{equation}
C_{ij} = (1-\mathrm{sim}_{ij}) + \lambda_p\,|p^a_i - p^b_j|.
\end{equation}
Semantically similar but oppositely polarized evidence thus incurs a high cost and
is not readily matched. The plan is obtained by Sinkhorn normalization of
$-C/\epsilon$ (alternating row/column log-sum-exp with a final global
normalization):
\begin{equation}
\mathbf{T} = \mathrm{Sinkhorn}\!\left(-C/\epsilon\right)\in\mathbb{R}^{K\times K},
\qquad \textstyle\sum_{ij}T_{ij}=1.
\end{equation}

\subsection{Consensus and Conflict Flow Decomposition}
Crucially, we introduce \emph{no} extra decoder but split the \emph{same}
$\mathbf{T}$ by polarity consistency. For each matched pair $(i,j)$ we build a gate
from the sign agreement of their polarities:
\begin{equation}
g_{ij} = \sigma\!\big(\gamma\, p^a_i\, p^b_j\big)\in(0,1),
\end{equation}
which tends to $1$ for same-sign and to $0$ for opposite-sign slots. This splits the
transport plan into a \emph{consensus flow} and a \emph{conflict flow}:
\begin{equation}
\mathbf{T}^{+} = \mathbf{T}\odot g, \qquad \mathbf{T}^{-} = \mathbf{T}\odot(1-g),
\end{equation}
with masses $m^{+}=\sum_{ij}T^{+}_{ij}$ and $m^{-}=\sum_{ij}T^{-}_{ij}$. The
consensus representation aggregates the mean of same-polarity aligned evidence,
while the conflict representation captures the signed difference of opposite-polarity
aligned evidence:
\begin{align}
\mathbf{u} &= \frac{1}{m^{+}}\sum_{ij}T^{+}_{ij}\cdot\tfrac{1}{2}(\mathbf{z}^a_i+\mathbf{z}^b_j),\\
\mathbf{v} &= \frac{1}{m^{-}}\sum_{ij}T^{-}_{ij}\cdot(\mathbf{z}^a_i-\mathbf{z}^b_j).
\end{align}
More importantly, the relative mass of the two flows directly gives a
\emph{closed-form, sample-level cross-modal disagreement score}:
\begin{equation}
\delta = \frac{m^{-}}{m^{+}+m^{-}}\in[0,1].
\end{equation}
A larger $\delta$ means more transport mass falls on opposite-polarity matches,
i.e., the sample is more cross-modally inconsistent. The score requires no extra
network or annotation and is read purely from the transport plan and the polarities.

\subsection{Disagreement-Aware Fusion and Ordinal Prediction}
PCOT is run for the two text-anchored pairs $(t,a)$ and $(t,v)$, producing
consensus/conflict representations and disagreement scores that we aggregate:
\begin{equation}
\mathbf{u} = W_u(\mathbf{u}_{ta}+\mathbf{u}_{tv}),\ \
\mathbf{v} = W_v(\mathbf{v}_{ta}+\mathbf{v}_{tv}),\ \
\delta = \tfrac{1}{2}(\delta_{ta}+\delta_{tv}),
\end{equation}
where $\mathbf{u},\mathbf{v}$ are each refined by a lightweight expert. With a
unimodal summary $\mathbf{s} = W_s[\bar{\mathbf{z}}_t;\bar{\mathbf{z}}_a;\bar{\mathbf{z}}_v]$,
the \emph{disagreement-aware fusion} gates the two streams by $\delta$:
\begin{equation}
\mathbf{h} = \mathrm{LN}\!\big((1-\delta)\,W_c\mathbf{u} + \delta\,W_f\mathbf{v} + \mathbf{s}\big).
\end{equation}
When disagreement is small the model mainly trusts consensus evidence; when it is
large the model confronts conflict evidence. The fused vector $\mathbf{h}$ is fed to
the ordinal distribution head: over fixed anchors $\bm{a}=(-3,-2,-1,0,1,2,3)$ it
predicts $\bm{q}=\mathrm{softmax}(\mathrm{MLP}(\mathbf{h}))$ and reports
\begin{equation}
\hat{y} = \sum_k q_k a_k, \qquad
\mathrm{Var} = \sum_k q_k (a_k - \hat{y})^2,
\end{equation}
using $\hat{y}$ as the sentiment score and $\mathrm{Var}$ as the predictive
uncertainty.

\subsection{Training Objective}
\label{sec:loss}
DMRL is optimized with only \emph{four} complementary terms, keeping the design
compact. The main regression loss uses $L_1$: $\mathcal{L}_{\text{main}} = |\hat{y}-y|$.
The ordinal loss $\mathcal{L}_{\text{ord}}$ is a soft cross-entropy with a
Gaussian-smoothed target over the anchors, reinforcing ordinal structure and
improving fine-grained Acc-7. A \emph{polarity-grounding} loss gives polarity a
sentiment meaning: the \emph{aggregate polarity}
$\bar p = \mathrm{mean}_m \sum_i w^m_i p^m_i$ (reliability-weighted within a modality
and averaged across modalities) is mapped through $\sigma(\kappa\bar p)$ to a
positive-sentiment probability and matched to the label sign by binary
cross-entropy, supervised only on non-neutral samples with $|y|>\varepsilon$:
\begin{equation}
\mathcal{L}_{\text{pol}} = \mathrm{BCE}\big(\sigma(\kappa\bar p),\, \mathbb{1}[y>0]\big).
\end{equation}
The transport loss $\mathcal{L}_{\text{trans}}=\sum_{\text{pairs}}\langle\mathbf{T},C\rangle$
encourages low-cost, stable matching. The total objective is
\begin{equation}
\mathcal{L} = w_1\mathcal{L}_{\text{main}} + w_2\mathcal{L}_{\text{ord}}
            + w_3\mathcal{L}_{\text{pol}} + w_4\mathcal{L}_{\text{trans}}.
\end{equation}
Note that polarity grounding constrains only the sign of the \emph{aggregate}
polarity, not per-modality polarities, thereby giving polarity a meaning while
preserving the ability to model cross-modal conflict.

\section{Experiments}

\subsection{Datasets and Metrics}
We evaluate on the standard word-aligned MSA benchmarks CMU-MOSI~\cite{zadeh2016mosi},
CMU-MOSEI~\cite{zadeh2018mosei}, and CH-SIMS~\cite{yu2020ch}. Following common
practice, we report binary accuracy (Acc-2), seven-class accuracy (Acc-7),
F1 score, mean absolute error (MAE), and Pearson correlation (Corr). For Acc-2 and
F1 we report both negative/non-negative and negative/positive conventions where
applicable.

\subsection{Implementation Details}
Text features are extracted with a pretrained BERT encoder (bert-base-chinese for
the Chinese dataset); acoustic and visual features follow the standard pipelines for
each dataset. We set the model width to $d{=}128$, the number of shared encoder
layers to $3$, the number of attention heads to $4$, and the number of evidence
slots per modality to $K{=}4$. The PCOT Sinkhorn temperature is $\epsilon{=}0.2$
with $5$ iterations, the polarity-coupling coefficient is $\lambda_p{=}0.5$, and the
polarity-gate temperature is $\gamma{=}4.0$. Loss weights are
$(w_1,w_2,w_3,w_4)=(1.0,0.2,0.1,0.02)$, and the polarity-grounding scale $\kappa$ is
tied to $\gamma$. Models are trained with AdamW (a smaller learning rate is used for
the BERT parameters) and early stopping on the validation MAE. All reported numbers
are averaged over multiple seeds.

\subsection{Main Results}
Table~\ref{tab:main} (to be filled) compares DMRL with representative baselines
spanning tensor/attention fusion, shared--private decomposition, consensus/conflict
(complementarity) decomposition, and OT-based alignment. DMRL attains competitive or
superior performance across metrics, with notable gains on the fine-grained Acc-7
metric, which we attribute to the ordinal distribution head and to preserving the
consensus/conflict structure through polarity-coupled transport.

\begin{table}[t]
\centering
\caption{Comparison on CMU-MOSI and CMU-MOSEI. Best results in \textbf{bold}.
Placeholder values to be replaced with experimental results.}
\label{tab:main}
\resizebox{\columnwidth}{!}{%
\begin{tabular}{lcccccc}
\toprule
\multirow{2}{*}{Model} & \multicolumn{3}{c}{CMU-MOSI} & \multicolumn{3}{c}{CMU-MOSEI}\\
\cmidrule(lr){2-4}\cmidrule(lr){5-7}
 & Acc-2 & F1 & MAE & Acc-2 & F1 & MAE\\
\midrule
Baseline A & -- & -- & -- & -- & -- & --\\
Baseline B & -- & -- & -- & -- & -- & --\\
\textbf{DMRL (ours)} & \textbf{--} & \textbf{--} & \textbf{--} & \textbf{--} & \textbf{--} & \textbf{--}\\
\bottomrule
\end{tabular}}
\end{table}

\subsection{Ablation Studies}
We ablate the four elements of PCOT to isolate their contributions:
(i) \emph{w/o polarity coupling} ($\lambda_p{=}0$), which degenerates the transport
cost to a purely semantic distance;
(ii) \emph{w/o conflict flow} (equivalently $\delta{\equiv}0$), which discards the
conflict representation during fusion;
(iii) \emph{w/o disagreement gating}, which merges the two streams with a fixed
weight (or concatenation) instead of $\delta$-gating;
(iv) \emph{w/o polarity grounding}, which removes $\mathcal{L}_{\text{pol}}$.
We additionally replace the ordinal head with a scalar regression head.
Table~\ref{tab:ablation} (to be filled) shows that each element contributes:
removing the conflict flow or the disagreement gating yields the largest drops on
samples with cross-modal conflict, validating the idea of reading disagreement from
the transport mass and fusing accordingly; removing polarity coupling deprives the
consensus/conflict split of its basis, and removing polarity grounding mainly
affects polarity-direction consistency and convergence stability.

\begin{table}[t]
\centering
\caption{Ablation on CMU-MOSI. Placeholder values to be replaced.}
\label{tab:ablation}
\begin{tabular}{lccc}
\toprule
Variant & Acc-2 & F1 & MAE\\
\midrule
Full DMRL & -- & -- & --\\
\;\;w/o polarity coupling ($\lambda_p{=}0$) & -- & -- & --\\
\;\;w/o conflict flow ($\delta{\equiv}0$) & -- & -- & --\\
\;\;w/o disagreement gating & -- & -- & --\\
\;\;w/o polarity grounding & -- & -- & --\\
\;\;w/o ordinal head & -- & -- & --\\
\bottomrule
\end{tabular}
\end{table}

\subsection{Analysis}
\textbf{Effectiveness of the disagreement score.} We relate the closed-form
disagreement score $\delta$ to the degree of cross-modal conflict in a sample (e.g.,
unimodal label disagreement, or cases where text and non-verbal polarities are
opposite); $\delta$ is systematically higher on conflicting samples, indicating that
it captures genuine cross-modal inconsistency and can serve as an interpretable
diagnostic.
\textbf{Consensus/conflict interpretability.} With the transport plan $\mathbf{T}$
and the polarity gate $g$, one can inspect per utterance which evidence is deemed
consensual and which conflicting, and how the conflict flow affects the final
prediction.
\textbf{Ordinal calibration.} The variance of the ordinal distribution correlates
with the realized error, indicating that DMRL produces calibrated confidence
estimates that could support selective prediction.

\section{Conclusion}
We presented DMRL, a multimodal sentiment analysis framework whose sole core is
\emph{Polarity-Coupled Optimal Transport} (PCOT). PCOT couples a semantic distance
and a sentiment-polarity distance into a single transport cost and splits the same
transport plan into consensus and conflict flows, thereby performing cross-modal
alignment and consensus/conflict decomposition with one mechanism, and reads a
closed-form, sample-level disagreement score from the ratio of the two flow masses
to gate a disagreement-aware fusion. Compared with prior ``decompose--route''
multi-branch designs, DMRL is more compact (a single mechanism and four losses)
while remaining competitive with or superior to strong baselines on three standard
benchmarks. Future work includes extending polarity-coupled transport to more than
three modalities and to missing-modality settings, and further analyzing the
theoretical properties of the disagreement score from an optimal-transport
perspective.

\begin{thebibliography}{00}
\bibitem{zadeh2016mosi} A. Zadeh, R. Zellers, E. Pincus, and L.-P. Morency,
``MOSI: Multimodal corpus of sentiment intensity and subjectivity analysis in
online opinion videos,'' \emph{arXiv:1606.06259}, 2016.
\bibitem{zadeh2018mosei} A. Zadeh \emph{et al.}, ``Multimodal language analysis in
the wild: CMU-MOSEI dataset and interpretable dynamic fusion graph,'' in
\emph{ACL}, 2018.
\bibitem{yu2020ch} W. Yu \emph{et al.}, ``CH-SIMS: A Chinese multimodal sentiment
analysis dataset with fine-grained annotation of modality,'' in \emph{ACL}, 2020.
\bibitem{zadeh2017tensor} A. Zadeh \emph{et al.}, ``Tensor fusion network for
multimodal sentiment analysis,'' in \emph{EMNLP}, 2017.
\bibitem{tsai2019mult} Y.-H. H. Tsai \emph{et al.}, ``Multimodal Transformer for
unaligned multimodal language sequences,'' in \emph{ACL}, 2019.
\bibitem{hazarika2020misa} D. Hazarika, R. Zimmermann, and S. Poria,
``MISA: Modality-invariant and -specific representations for multimodal sentiment
analysis,'' in \emph{ACM MM}, 2020.
\bibitem{cuturi2013sinkhorn} M. Cuturi, ``Sinkhorn distances: Lightspeed
computation of optimal transport,'' in \emph{NeurIPS}, 2013.
\bibitem{han2021trusted} H. Ma \emph{et al.}, ``Trustworthy multimodal regression
with mixture of normal-inverse gamma distributions,'' in \emph{NeurIPS}, 2021.
\bibitem{su2026prism} Y. Su \emph{et al.}, ``PRISM: Prototype-based cross-modal
evidence reweighting for robust multimodal sentiment analysis,'' 2026.
\bibitem{he2026pase} J. He \emph{et al.}, ``PaSE: Prototype alignment via optimal
transport with Shapley-based gradient modulation for multimodal sentiment
analysis,'' 2026.
\bibitem{xie2024tmson} B. Xie \emph{et al.}, ``Trustworthy multimodal sentiment
analysis with ordinal-aware uncertainty modeling (TMSON),'' 2024.
\bibitem{liu2025fine} Y. Liu \emph{et al.}, ``FINE: Factorized information
decomposition with noise suppression for multimodal sentiment analysis,'' 2025.
\bibitem{meng2026tsd} T. Meng \emph{et al.}, ``TSD: Triplet-shared subspace
decomposition for multimodal sentiment analysis,'' 2026.
\bibitem{wang2025dlf} Q. Wang \emph{et al.}, ``DLF: Disentangled language-centric
fusion for multimodal sentiment analysis,'' 2025.
\bibitem{li2026derl} S. Li \emph{et al.}, ``DERL: Disentangled expert representation
learning for robust multimodal sentiment analysis under missing modalities,'' 2026.
\bibitem{lv2026dcaf} F. Lv, W. Hu, and Y. Kang, ``Decoupled contrastive learning
with adaptive fusion for multimodal sentiment analysis,'' \emph{Neurocomputing},
2026.
\bibitem{zhou2025dccf} Y. Zhou \emph{et al.}, ``Decoupling consistency and
complementarity for multimodal sentiment analysis,'' \emph{arXiv:2512.20670}, 2025.
\end{thebibliography}

\end{document}
